\section{B.4. Exact Kummer images and the Fermat fibres}

For $a\in\mathbb Q^\times$ let
$K_N(a)=\{u\in\overline{\mathbb Q}:u^N=a\}$, with transition map
$u\mapsto u^{M/N}$ for $N\mid M$, and let
$K(a)=\varprojlim_NK_N(a)$. We fix the compatible positive roots
$\theta_N=2^{1/N}$ and roots of unity $\zeta_N=e^{2\pi i/N}$.
Let $\Gamma_{\mathbb Q}=\operatorname{Gal}(\overline{\mathbb Q}/\mathbb Q)$.
The conventions in cote 138, scans 35--41, then give
\[
 \sigma(\zeta_N)=\zeta_N^{a_N(\sigma)},\qquad
 \sigma(\theta_N)=\zeta_N^{b_N(\sigma)}\theta_N.
\]
Thus $\sigma$ acts as $x\mapsto ax+b$ on $K(2)$, where
$a\in\widehat{\mathbb Z}^{\times}$ and $b\in\widehat{\mathbb Z}$.
We write affine pairs in this order; their product is
$(a,b)(a',b')=(aa',b+ab')$.

\paragraph{B4.1 (Proved: the complete image for $2$ and $1/2$).}
\label{B4.1}
Define
\[
 \chi_2(a)=(-1)^{(a^2-1)/8}
 =\begin{cases}1,&a\equiv1,7\pmod8,\\-1,&a\equiv3,5\pmod8.
 \end{cases}
\]
The image of $\Gamma_{\mathbb Q}$ on $K(2)$ is exactly
\[
 H_2=\{(a,b):(-1)^b=\chi_2(a)\}
 \subset\widehat{\mathbb Z}^{\times}\ltimes\widehat{\mathbb Z}.
\]
It has index $2$. At every positive level $N$, its image has order
\[
 \frac{N\varphi(N)}{d_N},\qquad
 d_N=\begin{cases}2,&8\mid N,\\1,&8\nmid N.\end{cases}
\]
When $8\mid N$, the finite image is defined by the same parity
condition; otherwise it is the whole affine group.
With the reciprocal origin $\theta_N^{-1}$, the action on $K(1/2)$ is
$(a,b):x\mapsto ax-b$, and its image and indices are the same after
the change $b\mapsto-b$.

\begin{proof}
Put $F_N=\mathbb Q(\zeta_N)$ and
$L_N=F_N(\theta_N)$, the splitting field of $t^N-2$.
Every element of $\operatorname{Gal}(L_N/F_N)$ is determined by its
ratio on $\theta_N$, giving an injection into $\mu_N$. Let its order
be $d\mid N$. Its image is the unique subgroup $\mu_d$; therefore
$\theta_N^d=2^{1/(N/d)}$ lies in $F_N$.
Set $s=N/d$. The polynomial $t^s-2$ is Eisenstein at $2$ and hence
$\mathbb Q(2^{1/s})$ has degree $s$. Being a subfield of the abelian
extension $F_N/\mathbb Q$, it must be Galois. It is a real field,
whereas for $s>2$ the irreducible polynomial $t^s-2$ has nonreal roots.
Consequently $s\leq2$. It follows that $d=N$, except possibly
$d=N/2$ when $N$ is even. The latter alternative holds exactly when
$\sqrt2\in F_N$: necessity follows from the fixed radical, and
sufficiency follows because $\theta_N^{N/2}=\sqrt2$ restricts the
translation subgroup to even exponents.

We verify the cyclotomic condition without assuming a conductor formula.
The compositum of $F_N$ and $\mathbb Q(\zeta_8)$ is
$\mathbb Q(\zeta_{\operatorname{lcm}(N,8)})$, so their intersection
has degree
\[
 \frac{4\varphi(N)}{\varphi(\operatorname{lcm}(N,8))}.
\]
If $v_2(N)=0$ or $1$, this degree is $1$. If $v_2(N)=2$, it is $2$
and the intersection is $\mathbb Q(i)$, already contained in both
fields. None of these intersections contains $\sqrt2$. If $8\mid N$
the containment is immediate. This proves
$[L_N:F_N]=N/d_N$ and hence the asserted total order.

At any level divisible by $8$, the equality
$\sqrt2=\zeta_8+\zeta_8^{-1}$ gives
\[
 (-1)^{b_N(\sigma)}
 =\frac{\sigma(\sqrt2)}{\sqrt2}
 =\chi_2(a_N(\sigma)).
\]
This cuts out a subgroup of index $2$ in the finite affine group.
The already computed order proves equality. At the other levels the
computed order is the order of the entire affine group, so equality
holds there too. The profinite image is closed, since the absolute
Galois group is compact. It equals the inverse limit of these finite
images, which is exactly the displayed $H_2$. The sign constraint
cuts its ambient group into two cosets. Reciprocation of the roots
gives the assertion for $K(1/2)$.
\end{proof}

\paragraph{B4.2 (Proved: the complete image for $-1/2$).}\label{B4.2}
Take the compatible origins
$r_N=\zeta_{2N}\theta_N^{-1}\in K_N(-1/2)$ and put
$t(a)=(a-1)/2\in\widehat{\mathbb Z}$. This division is well-defined:
$a-1\in2\widehat{\mathbb Z}$ and multiplication by $2$ on
$\widehat{\mathbb Z}$ is injective. The translation coordinate on
$K(-1/2)$ is
\[
 c=t(a)-b,
\]
and the complete image is
\[
 H_{-1/2}=\{(a,c):(-1)^c=\chi_{-2}(a)\},\qquad
 \chi_{-2}(a)=(-1)^{(a-1)/2}\chi_2(a).
\]
Here $\chi_{-2}$ is $+1$ for $a\equiv1,3\pmod8$ and $-1$ for
$a\equiv5,7\pmod8$. Its adelic index is $2$ and its level-$N$ index
is again $d_N$.

\begin{proof}
At level $N$, one must know $a$ modulo $2N$ to compute the action on
$\zeta_{2N}$. Directly,
\[
 \frac{\sigma(r_N)}{r_N}
 =\zeta_N^{(a_{2N}(\sigma)-1)/2-b_N(\sigma)}.
\]
The identity $t(aa')=t(a)+a\,t(a')$ shows that
$(a,b)\mapsto(a,t(a)-b)$ is an involutive isomorphism of the profinite
affine group. Applying it to B4.1 gives the stated parity constraint.
When $8\mid N$, that constraint depends only on the level-$N$
coordinates and removes half the affine pairs. When $N$ is odd, a lift
of a translation modulo $N$ may have either parity. When $v_2(N)=1$
or $2$, a fixed unit modulo $N$ has lifts modulo $8$ with both values
of $\chi_{-2}$; the Chinese remainder theorem makes these choices
compatible with its odd components. Thus the projection is the full
affine group in all remaining cases. This proves every level index.
\end{proof}

\paragraph{B4.3 (Proved: all three Fermat fibres and all their indices).}
\label{B4.3}
Let $I=\{0,1,\infty\}$. In projective coordinates put
\[
 \overline C_N:\quad u_0^N+u_1^N+u_\infty^N=0,
 \qquad
 \overline C_N\longrightarrow\overline C_1,
 \quad [u_j]\longmapsto[u_j^N].
\]
On $\overline C_1$ use $z=-u_0/u_\infty$. The three points
$\alpha'(0)=2$, $\alpha'(1)=-1$, $\alpha'(\infty)=1/2$ have coordinate
$i$ equal to $1$ and the other two coordinates equal to $-1/2$.
For every $i$ their fibres satisfy
\[
 E_N(\alpha'(i))\cong K_N(-1/2)^2
\]
as $\Gamma_{\mathbb Q}$-sets and as torsors under
$\mu_N^I/\mu_N\cong\mu_N^2$.
With equal origins $(r_N,r_N)$, their complete image is
\[
 \{(a,(c,c)):(a,c)\in H_{-1/2}\}.
\]
With the opposite-phase origin used on scan~41, namely
$(\zeta_{2N}^{-1}\theta_N^{-1},\zeta_{2N}\theta_N^{-1})$,
it is
\[
 \left\{(a,(c-a+1,c)):(a,c)\in H_{-1/2}\right\}.
\]
Consequently the level-$N$ image has order $N\varphi(N)/d_N$ and its
index in the scalar affine group
$(\mathbb Z/N)^2\rtimes(\mathbb Z/N)^\times$ is $Nd_N$.
Its index in the full affine group with linear part
$\operatorname{GL}_2(\mathbb Z/N)$ is
\[
 Nd_N\frac{|\operatorname{GL}_2(\mathbb Z/N)|}{\varphi(N)}
 =d_NN^4\prod_{\ell\mid N}(1-\ell^{-2}).
\]
Both adelic ambient groups have infinite index. The same formulas hold
for each $i$, since permuting the three projective coordinates is
rational over $\mathbb Q$.

\begin{proof}
On the fibre over $\alpha'(i)$, the coordinate $u_i$ is nonzero.
Normalize it to $1$. The other two coordinates then have $N$th power
$-1/2$, which proves the asserted equivariant torsor isomorphism and
checks the sign in scan~35 directly from the Fermat equation.
The Galois automorphism acts on each coordinate by the same affine
pair $(a,c)$ of B4.2. This action is faithful as an action of that
finite affine group on the product: its restriction to varying either
coordinate determines $a$ and $c$. Therefore its image has the order
already computed in B4.2. The first coordinate of the opposite-phase
origin differs from $r_N$ by $\zeta_N^{-1}$. A change of origin by
$-1$ changes its translation to $c-(a-1)$, proving the second formula.
Counting the two ambient groups gives the indices. For the explicit last expression, reduction modulo $\ell$ gives $|\operatorname{GL}_2(\mathbb Z/\ell^k)|=\ell^{4(k-1)}(\ell^2-1)(\ell^2-\ell)$, and the Chinese remainder theorem multiplies these orders. Since $Nd_N$ is
unbounded, the closed adelic subgroup cannot be open and has infinite
index. No rank-two independence hypothesis is satisfied here: the two
radicands are equal.
\end{proof}

For the requested edge cases the orders and scalar-affine indices are
\[
\begin{array}{c|rrrr}
 N&1&2&4&8\\\hline
 |\operatorname{im}K_N(2)|&1&2&8&16\\
  \left[\operatorname{Aff}_1(N):\operatorname{im}K_N(2)\right]&1&1&1&2\\
  \left[\operatorname{Aff}_{2,\mathrm{scalar}}(N):\operatorname{im}E_N\right]&1&2&4&16
\end{array}
\]
Thus L4's explicit assertion for $K(2)$ and $K(1/2)$ is true.
Extending the same index $2$ to the full two-dimensional Fermat fibre
is false: already level $8$ has scalar-affine index $16$. B4.3 is the
complete replacement. The supplied study's $\sqrt2$ obstruction was
correct but did not by itself prove that there were no further
obstructions; B4.1 supplies that proof.

\paragraph{B4.4 (Proved under stated hypotheses: adelic openness for independent elements).}
\label{B4.4}
Let $K$ be a number field, let
$\xi_1,\dots,\xi_r\in K^\times$ be multiplicatively independent
modulo the roots of unity of $K$, and choose compatible roots of each.
The Kummer action gives a homomorphism
\[
 \rho:\Gamma_K\longrightarrow
 \widehat{\mathbb Z}^{r}\rtimes\widehat{\mathbb Z}^{\times},
\]
where the linear action is the common cyclotomic scalar.
Its image is open. More precisely, let $S$ contain the finite primes
occurring in the principal ideals of the $\xi_j$. In the finite free
abelian group
$U_S=\mathcal O_{K,S}^{\times}/\mu(K)$, let $H$ be the lattice
generated by the classes of the $\xi_j$, let
$H^{\mathrm{sat}}=(H\otimes\mathbb Q)\cap U_S$, and choose a positive
integer $e$ annihilating $H^{\mathrm{sat}}/H$. Put
$D=e|\mu(K)|$. Then
\[
 D\widehat{\mathbb Z}^{r}
 \subseteq\rho(\Gamma_{K(\mu_\infty)})
 \subseteq\widehat{\mathbb Z}^{r}.
\]
In particular the total index is at most
$[\widehat{\mathbb Z}^{\times}:\chi(\Gamma_K)]D^r$.
For $r=0$ this is simply the cyclotomic openness statement.
For $r\geq2$ the image is not open in the larger group
$\widehat{\mathbb Z}^{r}\rtimes\operatorname{GL}_r(\widehat{\mathbb Z})$.
If the $\xi_j$ are not rational, the acting group must be
$\Gamma_K$ for a field containing them; a separate permutation action
must be specified before replacing it by $\Gamma_{\mathbb Q}$.

\begin{proof}
The $S$-unit theorem makes $U_S$ finite free; the exact form used is
Milne, \emph{Algebraic Number Theory}~\cite{MilneANT2020}, version~3.08, Theorem~5.11.
Independence identifies $H$ with $\mathbb Z^r$, and its saturation
quotient is finite, so $e$ exists.
Put $L=K(\mu_\infty)$. We use the following exact part of Schinzel's
\emph{Abelian binomials, power residues and exponential congruences},
Theorem~2, p.~245~\cite{Schinzel1977}: if the splitting field of $t^N-a$ over a
characteristic-zero field $K$ is abelian and
$w_N=|\mu_N(K)|$, then $a^{w_N}\in K^{\times N}$.
The source was opened as a scan; the exponent is $w_N$, not $N$.

Suppose an integer vector $v=(v_j)$ gives
$a=\prod_j\xi_j^{v_j}\in L^{\times N}$.
All roots of $t^N-a$ then lie in $L$ because $L$ contains $\mu_N$.
Its splitting field over $K$ is abelian, as $L/K$ is abelian.
Schinzel's theorem gives $u\in K^\times$ with
$u^N=a^{w_N}$. Valuations outside $S$ show that $u$ is an $S$-unit.
Writing the identity in the free lattice $U_S$ gives
\[
 N[u]=w_N\sum_jv_j[\xi_j].
\]
Thus $[u]\in H^{\mathrm{sat}}$ and $e[u]\in H$. Comparing
coordinates in the basis $[\xi_j]$ of $H$ shows that
$N$ divides $ew_Nv_j$ for every $j$. Since $w_N$ divides $|\mu(K)|$,
we have $N\mid Dv_j$.

For completeness, the elementary Kummer pairing over $L$ is
\[
 \operatorname{Gal}(L(\xi_1^{1/N},\ldots,\xi_r^{1/N})/L)
 \times (\mathbb Z/N)^r\longrightarrow\mu_N,
\quad
 (\sigma,v)\longmapsto
 \prod_j\left(\frac{\sigma(\xi_j^{1/N})}{\xi_j^{1/N}}\right)^{v_j}.
\]
Its right kernel is exactly the vectors for which
$\prod_j\xi_j^{v_j}\in L^{\times N}$: triviality of the character
means the corresponding product of chosen radicals is fixed by the
Galois group, hence lies in $L$, and the converse follows by dividing
by an $N$th root in $L$. The character pairing of the finite group
$(\mathbb Z/N)^r$ is perfect, so a subgroup equals the annihilator of
its annihilator. The preceding divisibility therefore implies that
the Galois image contains $D(\mathbb Z/N)^r$.
Passing to the inverse limit, using compactness and extension of
finite-level automorphisms, gives the displayed inclusion of
$D\widehat{\mathbb Z}^{r}$. Its index is $D^r$.

The cyclotomic image $\chi(\Gamma_K)$ is open: its fixed field in
$\mathbb Q(\mu_\infty)$ is the finite extension
$K\cap\mathbb Q(\mu_\infty)$ of $\mathbb Q$.
The index of a subgroup projecting onto $\chi(\Gamma_K)$ equals
the index of its translation kernel, times the index of that
cyclotomic projection in the full scalar group. This proves the bound.
For $r\geq2$, scalar matrices have infinite index in
$\operatorname{GL}_r(\widehat{\mathbb Z})$: their indices modulo
primes are unbounded. Projecting the claimed larger affine group to
its linear factor proves the last assertion.

We can also derive the openness conclusion directly from the opened
Perucca--Sgobba theorem~\cite{PeruccaSgobba2019}, Theorem~1 of the
opened author preprint.
For $G=\langle\xi_1,\ldots,\xi_r\rangle$, their result supplies a
constant $C$ such that
\[
 \frac{N^r}{[K(\zeta_N,G^{1/N}):K(\zeta_N)]}\mid C
 \quad\text{for every }N\geq1.
\]
Independence makes $G$ torsion-free of rank $r$, so the hypotheses hold.
The finite affine image has order $[K(\zeta_N,G^{1/N}):K]$.
Its index is consequently at most
$C[\widehat{\mathbb Z}^{\times}:\chi(\Gamma_K)]$, uniformly in $N$.
These bounds imply finite index of the closed profinite image:
any finite set of distinct cosets remains distinct in some finite
quotient, so there cannot be more cosets than the common bound.
This independently derives the requested adelic conclusion from an
opened theorem for precisely the split torus in question.
\end{proof}

This derivation gives the prime-level maximality for every prime not
dividing $D$, and the required uniform control at the remaining primes.
A statement only for all sufficiently large primes, by itself, would
not establish adelic openness. Bashmakov's opened 1972 article
\cite{Bashmakov1972}, \S5, Lemma following Theorem~5, p.~53, gives
a fixed-prime finite-index result for the elliptic-curve Tate module
under the endomorphism hypotheses introduced on pp.~50--53. It does
not state the present split-torus adelic assertion. Ribet's 1979
article was located through its author's archive, but the primary
text remained blocked after checking the current and legacy publisher
links; we do not attribute an unchecked theorem number to it.
The two derivations above use opened sources with the precise needed
hypotheses. The conclusion for the rational elements in scan~20
is obtained by taking $K=\mathbb Q$; its ambient affine group must
have the scalar linear part indicated in that scan.
